\section{Second Pass: Logic as a Language You Can Speak}
\chapterline

\subsection{Beyond the readiness plan: from templates to spoken dialect}

In Chapter~26, we used explicit proof templates for direct arguments, contrapositives, and contradictions to get clean arguments onto the page.

In this Second Pass, we focus on fluency. At a blackboard, mathematicians do not pause to consult templates or compute sixteen-row truth tables. They read and evaluate logic directly, tracking how truth values flow through a compound statement.

Efforts to formalize logical deduction gathered momentum in the late nineteenth century. In 1879, the German mathematician and philosopher Gottlob Frege published \emph{Begriffsschrift} (``Concept-Script''). Seeking to eliminate conversational ambiguity and unstated assumptions from mathematical proofs, Frege aimed to realize Gottfried Wilhelm Leibniz's earlier vision of a formal symbolic language in which every inference is governed by explicit rules.

In his two-volume work, \emph{Grundgesetze der Arithmetik} (``Basic Laws of Arithmetic''), Frege attempted to derive arithmetic from logical axioms. Yet when students first encounter formal logic, truth tables often feel like tedious accounting: listing mechanical rows of $\mathbf{T}$'s and $\mathbf{F}$'s, checking column by column whether two expressions match.

In practice, mathematicians treat compound logical statements as syntax trees: branched structures that break a composite sentence into its core connectives and sub-clauses, showing which operation evaluates first and which governs the whole statement. They see nested operators passing truth values upward from premises to conclusions, just like logic gates in a circuit.

\begin{center}
\begin{tikzpicture}[
  level distance=1.3cm,
  sibling distance=2.8cm,
  >=Latex,
  gate/.style={draw=structurepurple!80!black, fill=structurepurple!15, thick, rounded corners=3pt, font=\small\bfseries, inner sep=4pt},
  leaf/.style={draw=sourceblue!80!black, fill=sourceblue!15, thick, rounded corners=3pt, font=\small\bfseries, inner sep=4pt},
  value/.style={font=\scriptsize\bfseries, text=trapred}
]

  % Root: Implication
  \node[gate] (root) {$\implies$}
    child {
      node[gate] (and) {$\land$}
        child {
          node[leaf] (p) {$P$}
          node[below=2pt of p, font=\scriptsize, text=targetgreen!80!black] {$\mathbf{T}$}
        }
        child {
          node[gate, xshift=-0.5cm] (not) {$\neg$}
            child {
              node[leaf] (q) {$Q$}
              node[below=2pt of q, font=\scriptsize, text=trapred] {$\mathbf{F}$}
            }
          node[above right=1pt of not, font=\scriptsize, text=targetgreen!80!black] {$\mathbf{T}$}
        }
      node[above left=1pt of and, font=\scriptsize, text=targetgreen!80!black] {$\mathbf{T}$}
    }
    child {
      node[leaf, xshift=0.5cm] (r) {$R$}
      node[below=2pt of r, font=\scriptsize, text=trapred] {$\mathbf{F}$}
    };

  % Annotation for the evaluated output
  \node[above=4pt of root, font=\small\bfseries, text=trapred] {Result: $\mathbf{T} \implies \mathbf{F} = \mathbf{F}$};

  % Explanatory labels on the side
  \node[right=2.2cm of root, align=left, font=\small] {
    \textbf{Semantic Signal Flow}: \\
    1. Input: $P=\mathbf{T}, \, Q=\mathbf{F}, \, R=\mathbf{F}$ \\
    2. Sub-tree: $\neg Q = \neg \mathbf{F} = \mathbf{T}$ \\
    3. Conjunction: $P \land \neg Q = \mathbf{T} \land \mathbf{T} = \mathbf{T}$ \\
    4. Implication: $\mathbf{T} \implies \mathbf{F} = \mathbf{F}$
  };

\end{tikzpicture}
\end{center}

\subsection{Translating subtle English idioms into boolean structure}

Ordinary speech is rich with colloquial phrasing that obscures logical meaning. To speak logic fluently, you must recognize standard idiomatic translations:

\begin{itemize}
\item \textbf{``$A$ unless $B$''}: This means that if $B$ does not occur, $A$ must hold:
\[
\neg B \implies A \quad (\text{or equivalently, } A \lor B).
\]
For example, ``the sequence converges unless it oscillates'' means: if the sequence does not oscillate, it must converge.

\item \textbf{``$P$ if $Q$'' versus ``$P$ only if $Q$''}:
\begin{itemize}
\item ``$P$ if $Q$'' means $Q \implies P$. Here $Q$ is the hypothesis: whenever $Q$ holds, $P$ is guaranteed.
\item ``$P$ only if $Q$'' means $P \implies Q$. Notice the reversed direction! Saying ``you can pass only if you take the exam'' means that passing cannot happen without the exam ($\text{Pass} \implies \text{Exam}$). It does not say that taking the exam guarantees passing.
\item ``$P$ if and only if $Q$'' means both implications hold simultaneously: $(P \implies Q) \land (Q \implies P)$, written compactly as $P \iff Q$.
\end{itemize}

\item \textbf{Sufficient versus Necessary}: In any implication $P \implies Q$:
\begin{itemize}
\item $P$ is a \textbf{sufficient condition} for $Q$: having $P$ is enough on its own to guarantee $Q$.
\item $Q$ is a \textbf{necessary condition} for $P$: without $Q$, you cannot have $P$.
\end{itemize}
For a concrete memory hook: living in Kyoto is sufficient for living in Japan; living in Japan is necessary for living in Kyoto.

\begin{center}
\begin{tikzpicture}[scale=0.92, >=Latex]
  % Left panel: Euler Diagram of Inclusion P \subseteq Q
  \begin{scope}[shift={(-3.6, 0)}]
    % Outer set Q
    \draw[thick, fill=blue!10, draw=sourceblue!80!black] (0, 0) ellipse (2.3 and 1.4);
    \node[sourceblue!90!black, font=\scriptsize\bfseries] at (1.05, 0.15) {$Q$ (Japan)};
    \node[font=\tiny, sourceblue!80!black] at (1.05, -0.15) {(Necessary)};
    
    % Inner set P
    \draw[thick, fill=targetgreen!20, draw=targetgreen!80!black] (-0.95, 0) ellipse (0.85 and 0.7);
    \node[targetgreen!90!black, font=\scriptsize\bfseries] at (-0.95, 0.15) {$P$ (Kyoto)};
    \node[font=\tiny, targetgreen!80!black] at (-0.95, -0.15) {(Sufficient)};

    \node[font=\footnotesize, align=center, text=gray!80!black] at (0, -1.8) {
      $P \subseteq Q \iff (P \implies Q)$
    };
  \end{scope}

  % Right panel: Idiomatic Translation Table
  \begin{scope}[shift={(3.4, 0)}]
    \draw[rounded corners=4pt, fill=purple!3!white, draw=structurepurple!50, thick] (-3.1, -1.8) rectangle (3.1, 1.8);
    
    \node[font=\footnotesize\bfseries, structurepurple!90!black] at (0, 1.4) {Direction of Implication};

    \node[anchor=west, font=\scriptsize] at (-2.8, 0.7) {\textbf{``$P$ if $Q$''}: \hfill $Q \implies P$};
    \node[anchor=west, font=\scriptsize, text=gray!70!black] at (-2.6, 0.35) {($Q$ is hypothesis; $Q$ inside $P$)};

    \node[anchor=west, font=\scriptsize] at (-2.8, -0.1) {\textbf{``$P$ only if $Q$''}: \hfill $P \implies Q$};
    \node[anchor=west, font=\scriptsize, text=gray!70!black] at (-2.6, -0.45) {($Q$ is required; $P$ inside $Q$)};

    \node[anchor=west, font=\scriptsize] at (-2.8, -0.9) {\textbf{``$P$ iff $Q$''}: \hfill $P \iff Q$};
    \node[anchor=west, font=\scriptsize, text=gray!70!black] at (-2.6, -1.25) {($P$ and $Q$ coincide: $P = Q$)};
  \end{scope}
\end{tikzpicture}\\[0.3em]
{\small\textbf{The Geometry of Implication:} The statement $P \implies Q$ corresponds to set containment $P \subseteq Q$. Being inside Kyoto ($P$) guarantees being inside Japan ($Q$), making $P$ \emph{sufficient}. Being outside Japan ($\neg Q$) makes being in Kyoto impossible ($\neg P$), making $Q$ \emph{necessary}.}
\end{center}

\item \textbf{``Neither $P$ nor $Q$''}: Both conditions fail simultaneously:
\[
\neg P \land \neg Q \quad (\text{which, by de Morgan's laws, is } \neg(P \lor Q)).
\]

\item \textbf{``$P$ but $Q$''}: In everyday English, ``but'' conveys psychological surprise or contrast. In mathematics, it is logically identical to ``and'' ($P \land Q$). Both clauses are claimed to hold. For example, ``the function $f$ is continuous on $[0,1]$, but nowhere differentiable'' is the formal conjunction $\text{Continuous} \land \neg\text{Differentiable}$.
\end{itemize}

\begin{recoverybox}
\textbf{Pause and think}: A sign on a library reading room says: ``You may enter only if you display a faculty badge.'' A student observes Professor Klein displaying a valid badge, but the librarian does not let her in because the room is closed for renovations. Has the rule posted on the sign been violated?

\emph{Answer}: No! Displaying a faculty badge is a \emph{necessary} condition ($\text{Enter} \implies \text{Badge}$), not a \emph{sufficient} one. The rule guarantees that anyone who enters has a badge; it does not promise that having a badge forces the doors open.
\end{recoverybox}

\subsection{Short-circuit evaluation: why mathematicians skip truth tables}

Consider the proposition:
\[
(A \land B \land C) \implies D.
\]
A student approaching this mechanically might write out a truth table with $2^4 = 16$ rows. But mathematicians never write such tables. Why?

The reason lies in the truth conditions of implication. An implication $P \implies Q$ is false in exactly one case: when the hypothesis $P$ is true and the conclusion $Q$ is false. When the hypothesis $P$ is false, the implication is true by definition---a condition known as \textbf{vacuously true}.

Now inspect $(A \land B \land C) \implies D$. How many rows of the 16-row truth table have a false antecedent?
\begin{itemize}
\item If $A = \mathbf{F}$ (8 rows), the conjunction $(A \land B \land C)$ is false immediately.
\item If $A = \mathbf{T}$ but $B = \mathbf{F}$ (4 rows), the conjunction is false.
\item If $A = \mathbf{T}$ and $B = \mathbf{T}$ but $C = \mathbf{F}$ (2 rows), the conjunction is false.
\end{itemize}
Summing these gives $8 + 4 + 2 = 14$ rows where the premise is false! In all 14 of those rows, the implication evaluates to $\mathbf{T}$ automatically. No mathematical calculation is needed.

Only $16 - 14 = 2$ rows remain: the scenario where $A, B$, and $C$ are all true. In those two rows, the truth of the whole statement depends entirely on whether $D$ is true.

\begin{center}
\begin{tikzpicture}[
  >=Latex,
  font=\small,
  box/.style={draw=#1!80!black, fill=#1!10, thick, rounded corners=4pt, inner sep=5pt, align=center},
  pipe/.style={draw=sourceblue!80!black, fill=sourceblue!8, thick, rounded corners=4pt, inner sep=5pt, align=center}
]

  \node[font=\bfseries, text=structurepurple] at (-3.7, 2.7) {The 16 Truth-Table Worlds};

  \node[box=trapred, text width=5.0cm] (vac) at (-3.7, 0.95) {
    \textbf{14 Rows: Premise is False}\\[0.2em]
    \scriptsize $A=\mathbf{F}$ (8 rows) $\lor$ $B=\mathbf{F}$ (4 rows) $\lor$ $C=\mathbf{F}$ (2 rows)\\[0.3em]
    \footnotesize $\mathbf{F} \implies D \;\equiv\; \mathbf{T}$ \\
    \textbf{\textcolor{trapred!90!black}{Vacuously True by Default}}
  };

  \node[box=targetgreen, text width=5.0cm] (live) at (-3.7, -1.6) {
    \textbf{2 Rows: Premise is True}\\[0.2em]
    \scriptsize $A = \mathbf{T}, \; B = \mathbf{T}, \; C = \mathbf{T}$\\[0.2em]
    \footnotesize Truth depends entirely on $D$
  };

  \node[font=\bfseries, text=sourceblue] at (3.7, 2.7) {The Mathematician's Proof};

  \node[pipe, text width=5.2cm] (step1) at (3.7, 1.8) {
    \textbf{Step 1: Assume the Premise}\\[0.1em]
    \footnotesize ``Suppose $A, B$, and $C$ hold.''\\[0.1em]
    \scriptsize (Discards all 14 vacuous rows)
  };

  \node[pipe, text width=5.2cm] (step2) at (3.7, 0.1) {
    \textbf{Step 2: Deduce Target}\\[0.1em]
    \footnotesize Apply definitions, lemmas, and algebra\\
    \footnotesize under the active premise.
  };

  \node[pipe, text width=5.2cm, draw=targetgreen!80!black, fill=targetgreen!10] (step3) at (3.7, -1.6) {
    \textbf{Step 3: Establish Conclusion}\\[0.1em]
    \footnotesize Prove $D$ must hold.\\
    \scriptsize $\mathbf{T} \implies \mathbf{T} \;\equiv\; \mathbf{T}$. Done!
  };

  \draw[->, thick, sourceblue!80!black] (step1) -- (step2);
  \draw[->, thick, sourceblue!80!black] (step2) -- (step3);

  \draw[->, very thick, targetgreen!80!black] (live.east) -- (step3.west)
    node[midway, above=2pt, font=\scriptsize\bfseries, text=targetgreen!90!black] {Focus here};

\end{tikzpicture}

\vspace{0.3em}
\small\textbf{Figure 27.2}: The structure of direct proof for conditional statements. Rather than verifying 16 rows in a truth table, a direct proof assumes the antecedent, dismissing the 14 vacuous cases and focusing directly on establishing the conclusion $D$.
\end{center}

This short-circuiting explains why proofs of $P \implies Q$ invariably begin: \emph{``Assume $P$.''} You do not need to analyze the worlds where $P$ is false, because the statement is already true in all of them.

\subsection{Four lenses on a single mathematical claim}

To see how fluency transforms proof writing, consider this elementary claim:
\[
\forall n \in \Z, \; n^2 + n \text{ is even}.
\]
Let us examine this statement through four distinct perspectives: reading, writing, debugging, and structural algebra.

\subsubsection{1. The Reader's Lens: Unpacking grammar and target}
The outer symbol is universal: $\forall n \in \Z$. A valid proof must hold for an arbitrary, unspecified integer $n$. We cannot pick $n = 4$ and call it a day.
The property to prove is ``$n^2 + n$ is even.'' By definition, this means:
\[
\exists k \in \Z \text{ such that } n^2 + n = 2k.
\]
Our algebraic target is now concrete: we must factor or rearrange $n^2 + n$ until a factor of $2$ appears with an integer quotient.

\subsubsection{2. The Writer's Lens: Crafting clear, self-contained prose}
How do we guide the reader from the arbitrary integer to the factor of $2$?
\begin{proof}
Let $n \in \Z$ be arbitrary. We can factor the expression as:
\[
n^2 + n = n(n+1).
\]
The integers $n$ and $n+1$ are consecutive. By the division algorithm, every integer is either even or odd:
\begin{itemize}
\item \textbf{Case 1 ($n$ is even)}: There exists $m \in \Z$ such that $n = 2m$. Substituting this gives
\[
n(n+1) = 2m(n+1) = 2(m(n+1)).
\]
Since $m(n+1)$ is an integer, $n^2 + n$ is even.
\item \textbf{Case 2 ($n$ is odd)}: There exists $m \in \Z$ such that $n = 2m + 1$. Then
\[
n + 1 = (2m + 1) + 1 = 2m + 2 = 2(m+1).
\]
Substituting this factor into the product yields
\[
n(n+1) = n \cdot 2(m+1) = 2(n(m+1)).
\]
Since $n(m+1)$ is an integer, $n^2 + n$ is even.
\end{itemize}
In both cases, $n^2 + n$ is twice an integer, completing the proof.
\end{proof}

\subsubsection{3. The Debugger's Lens: Diagnosing flawed reasoning}
Consider this common student draft:
\begin{quote}
\emph{Attempted proof}: $n^2 + n$ is even because it has two terms, one of which is even and one of which is odd, and even plus odd is always odd, wait, no, they cancel out.
\end{quote}
\textbf{Diagnosis}:
\begin{itemize}
\item The claim that $n^2$ and $n$ have opposite parity is completely false! If $n = 2$, both $n = 2$ and $n^2 = 4$ are even. If $n = 3$, both $n = 3$ and $n^2 = 9$ are odd.
\item In fact, $n^2$ and $n$ \emph{always} have the exact same parity: an odd times an odd is odd, and an even times an even is even.
\item The sum of two numbers with identical parity is always even ($\text{odd} + \text{odd} = \text{even}$, and $\text{even} + \text{even} = \text{even}$).
\end{itemize}
Debugging the flawed intuition leads directly to an alternative, highly elegant proof: analyzing the parity of $n$ directly!

\subsubsection{4. The Structural Lens: Connecting to modular arithmetic}
In higher algebra and number theory, divisibility statements are expressed as congruences modulo $2$---clock arithmetic with only two marks on the dial, $0$ and $1$, where integers are grouped together if they leave the same remainder when divided by $2$:
\[
n^2 + n \equiv n(n+1) \pmod 2.
\]
In the two-element arithmetic system $\Z/2\Z = \{[0], [1]\}$ (often called the finite field $\mathbb{F}_2$, where addition and multiplication operate modulo $2$):
\begin{align*}
[0]([0] + [1]) &= [0] \cdot [1] = [0], \\
[1]([1] + [1]) &= [1] \cdot [0] = [0].
\end{align*}
The polynomial $x^2 + x$ evaluates to zero on every element of $\mathbb{F}_2$. What began as an algebraic calculation with consecutive integers is recognized as a fundamental property of polynomial evaluation over finite fields!

\begin{center}
\begin{tikzpicture}[scale=0.92, >=Latex]
  % Panel 1: Reader
  \node[draw=sourceblue!80!black, fill=sourceblue!8, rounded corners=4pt, thick, text width=5.6cm, align=left, font=\scriptsize, inner sep=5pt] at (-3.3, 1.3) {
    \textbf{\textcolor{sourceblue!90!black}{1. Reader's Lens: Quantifier Contract}}\\[0.2em]
    $\forall n \in \Z, \; \exists k \in \Z \text{ such that } n^2 + n = 2k$.\\[0.1em]
    Goal: Given arbitrary $n$, extract factor of $2$.
  };

  % Panel 2: Writer
  \node[draw=targetgreen!80!black, fill=targetgreen!8, rounded corners=4pt, thick, text width=5.6cm, align=left, font=\scriptsize, inner sep=5pt] at (3.3, 1.3) {
    \textbf{\textcolor{targetgreen!90!black}{2. Writer's Lens: Consecutive Parity}}\\[0.2em]
    Factor $n^2+n = n(n+1)$. By integer parity:\\[0.1em]
    $\bullet$ $n = 2m \implies 2[m(n+1)]$\\[0.1em]
    $\bullet$ $n = 2m+1 \implies 2[n(m+1)]$.
  };

  % Panel 3: Debugger
  \node[draw=quotientorange!80!black, fill=quotientorange!8, rounded corners=4pt, thick, text width=5.6cm, align=left, font=\scriptsize, inner sep=5pt] at (-3.3, -1.0) {
    \textbf{\textcolor{quotientorange!90!black}{3. Debugger's Lens: Shared Parity}}\\[0.2em]
    $n^2$ and $n$ always share identical parity:\\[0.1em]
    $\bullet$ $\text{even}^2 + \text{even} = \text{even} + \text{even} = \text{even}$\\[0.1em]
    $\bullet$ $\text{odd}^2 + \text{odd} = \text{odd} + \text{odd} = \text{even}$.
  };

  % Panel 4: Structural
  \node[draw=structurepurple!80!black, fill=structurepurple!8, rounded corners=4pt, thick, text width=5.6cm, align=left, font=\scriptsize, inner sep=5pt] at (3.3, -1.0) {
    \textbf{\textcolor{structurepurple!90!black}{4. Structural Lens: Finite Field $\mathbb{F}_2$}}\\[0.2em]
    Polynomial $x^2 + x = x(x+1)$ in $\Z/2\Z = \{0, 1\}$:\\[0.1em]
    $\bullet$ $x = 0 \implies 0 \cdot 1 \equiv 0 \pmod 2$\\[0.1em]
    $\bullet$ $x = 1 \implies 1 \cdot 0 \equiv 0 \pmod 2$.
  };
\end{tikzpicture}\\[0.3em]
{\small\textbf{Four Lenses on One Truth:} The simple claim $\forall n \in \Z, n^2+n \text{ is even}$ seen through the lenses of formal contract, constructive case breakdown, parity invariance, and finite-field roots.}
\end{center}

\begin{labbox}
Pick any proposition with two connectives, such as $(P \implies Q) \lor (Q \implies P)$. Draw its abstract syntax tree, evaluate it for all four truth assignments, and observe why it evaluates to true in every possible case (a classical tautology). Then test Peirce's law: $((P \implies Q) \implies P) \implies P$. Notice how the tree structure helps you isolate which branch forces truth even when you cannot determine whether $Q$ holds.
\end{labbox}

\begin{historybox}
The development of formal logic led to foundational debates in the early twentieth century. In June 1902, as the second volume of Gottlob Frege's \emph{Grundgesetze} was in press, Bertrand Russell pointed out that Frege's logical system permitted unrestricted comprehension---the assumption that any predicate defines a set. Russell introduced the collection $R = \{x : x \notin x\}$ of all sets that do not contain themselves. If $R \in R$, then by definition $R \notin R$; if $R \notin R$, then $R \in R$. This contradiction, now known as Russell's paradox, showed that naive set theory required reformulation.

To resolve these paradoxes, Russell and Alfred North Whitehead developed type theory in \emph{Principia Mathematica} (1910--1913), stratifying mathematical objects into a hierarchy of types to prevent a set from containing itself.

Alternative perspectives emerged. L.E.J. Brouwer founded \emph{intuitionism}, arguing that mathematical objects must be constructively created and questioning the unrestricted use of the law of excluded middle on infinite collections. In response, David Hilbert formulated the \emph{formalist program}, aiming to prove arithmetic to be both consistent (free of contradictions) and complete (capable of proving or refuting every statement).

In 1931, Kurt G\"odel proved his Incompleteness Theorems, demonstrating that any consistent, effectively generated formal system capable of expressing elementary arithmetic is incomplete: there exist true arithmetic statements that cannot be proven within the system. Nor can such a system prove its own consistency. In 1936, Alan Turing addressed Hilbert's \emph{Entscheidungsproblem} by formalizing the concept of an abstract computing machine and proving the undecidability of the Halting Problem, establishing foundational results for theoretical computer science.

Today, formal logic remains closely tied to computer science. Interactive proof assistants such as Lean, Coq, and Isabelle allow mathematicians to verify formal proofs mechanically, connecting classical proof theory with software verification.
\end{historybox}

\subsection{Exercises}

\textbf{Part A: Mental Checks and Syntax Parsing.} For each proposition, draw its abstract syntax tree and identify its root (outermost) connective:
\begin{enumerate}
\item $(P \land \neg Q) \implies (R \lor S)$.
\item $\neg(P \implies Q) \iff (P \land \neg Q)$.
\item $(P \implies Q) \land (Q \implies R) \implies (P \implies R)$.
\item Trace the truth signal flow for the tree in Exercise 1 under the truth assignment $P = \mathbf{T}, Q = \mathbf{T}, R = \mathbf{F}, S = \mathbf{F}$. What is the truth value at the root?
\item \emph{Mental check}: Suppose an implication $P \implies Q$ is known to evaluate to $\mathbf{F}$. What can you deduce about the truth value of $P$? What can you deduce about the truth value of $Q$?
\end{enumerate}

\textbf{Part B: Translating Spoken Idioms.} Translate each English statement into precise propositional formulas using the indicated letters:
\begin{enumerate}[resume]
\item ``The sequence $(a_n)$ is bounded unless it diverges to infinity'' ($B$: sequence is bounded, $D$: sequence diverges to infinity).
\item ``A quadrilateral is a rectangle only if its interior angles are right angles'' ($R$: quadrilateral is a rectangle, $A$: interior angles are right angles).
\item ``A function $f$ is integrable if it is continuous'' ($I$: function is integrable, $C$: function is continuous).
\item ``The linear system has a unique solution if and only if the matrix determinant is non-zero'' ($U$: unique solution, $D$: determinant is non-zero).
\item Identify the sufficient condition and the necessary condition in the theorem: ``Every differentiable function is continuous.''
\item Explain the logical difference between the mathematical assertion ``$f$ is continuous, but not bounded'' and the compound statement ``$f$ is continuous and not bounded.''
\end{enumerate}

\textbf{Part C: Short-Circuiting and Vacuous Truth.} Examine each implication and determine whether it is vacuously true, non-vacuously true, or false:
\begin{enumerate}[resume]
\item ``If $x$ is a positive real number such that $x^2 + 1 = 0$, then $x + 10 = 3$.''
\item ``If $n$ is an even prime integer greater than $2$, then $n$ is a multiple of $7$.''
\item In a textbook proof of Lagrange's Theorem (``If $G$ is a finite group and $H$ is a subgroup of $G$, then the order of $H$ divides the order of $G$''), describe which scenarios are dismissed as vacuously true by beginning the proof with ``Let $G$ be a finite group and let $H \le G$.''
\end{enumerate}

\textbf{Part D: Proofs Through the Four Lenses.} Apply the four lenses (reader, writer, debugger, structural) to each claim:
\begin{enumerate}[resume]
\item Prove that for every integer $n$, the product $n(n+1)(n+2)$ is divisible by $6$.
\begin{enumerate}
\item \emph{Reader's Lens}: State the explicit existential definition of divisibility by $6$.
\item \emph{Writer's Lens}: Write a clean, self-contained proof using properties of consecutive integers or division by cases.
\item \emph{Debugger's Lens}: Diagnose this student attempt: ``$n(n+1)$ is divisible by $2$, and $(n+1)(n+2)$ is divisible by $3$, so multiplying them together gives divisibility by $6$.''
\item \emph{Structural Lens}: Evaluate the polynomial $x(x+1)(x+2)$ in the modular rings $\Z/2\Z$ and $\Z/3\Z$.
\end{enumerate}
\item Prove that the square of any odd integer leaves a remainder of $1$ when divided by $8$. (Hint: Write $n = 2k + 1$ and consider the parity of the consecutive product $k(k+1)$).
\end{enumerate}

\textbf{Part E: Deeper Structural Investigations.}
\begin{enumerate}[resume]
\item \textbf{Peirce's Law}: Consider the classical propositional formula $((P \implies Q) \implies P) \implies P$.
\begin{enumerate}
\item Evaluate the formula when $P = \mathbf{T}$ using short-circuit reasoning.
\item Evaluate the formula when $P = \mathbf{F}$.
\item Why is this formula striking? It contains only implications ($\implies$), yet it asserts the law of excluded middle without using disjunction ($\lor$) or negation ($\neg$).
\end{enumerate}
\item \textbf{Functional Completeness of the Sheffer Stroke}: The NAND operator, written $P \mid Q$, is defined as $\neg(P \land Q)$.
\begin{enumerate}
\item Show that $\neg P \equiv P \mid P$.
\item Show that $P \lor Q \equiv (P \mid P) \mid (Q \mid Q)$.
\item Conclude that every propositional connective can be expressed using only the Sheffer stroke.
\end{enumerate}
\end{enumerate}

\begin{walkbox}
When reading your next mathematics paper or textbook chapter, notice the outermost connective of every theorem. Does the theorem say ``If $\dots$ then'', ``$\dots$ if and only if $\dots$'', or ``Neither $\dots$ nor $\dots$''? Notice how identifying the root connective in your head immediately tells you what the first sentence of the proof must be.
\end{walkbox}

\subsection{The limits of propositional truth---and the quantifier shift}

Once you read logical connectives as grammar rather than decoration, the shape of a statement tells you how to handle it. Finding the main connective in the syntax tree shows how a proof must open; seeing why $P \implies Q$ is vacuously true when $P$ fails explains why a direct proof starts by assuming $P$ and working toward $Q$; and switching among the viewpoints of reader, writer, debugger, and Boolean algebraist lets you both build arguments and negate them cleanly.

Yet propositional logic operates on whole statements. It treats propositions $P$ and $Q$ as indivisible units, without analyzing the individual elements and predicates within them. It cannot express statements involving quantification over elements, such as ``every integer has a prime factorization'' or ``there exists a unique solution.''

To analyze statements where elements and quantifiers interact, Chapter~28 presents \textbf{Second Pass: Quantifier Order in Real Problems}, examining how the order of $\forall$ and $\exists$ distinguishes pointwise from uniform convergence.

\newpage
